\documentclass[11pt]{article}
\usepackage[margin=1in]{geometry}
\usepackage[T1]{fontenc}
\usepackage{lmodern}
\usepackage{amsmath,amssymb,amsthm}
\usepackage[hidelinks]{hyperref}
\usepackage{microtype}
\raggedright

\newtheorem*{theorem}{Theorem}
\newtheorem*{proposition}{Proposition}
\newcommand{\D}{\mathbb D}
\newcommand{\B}{\mathbb B}

\begin{document}

\begin{center}
{\Large Literature-implied full answer for arXiv:math/0702269}\\[4pt]
{\large Balanced monomials violate the Schur--Agler ball estimate}
\end{center}

\paragraph{Status.}
\textbf{literature\_implied\_answer (full negative answer).}
The counterexample below completely answers the existence question in Open
Problem~2 of Anderson--Dritschel--Rovnyak.  Its ingredients are already implicit
in the sharp coefficient theorem of Dai--Chen--Pan, but the supporting paper
does not cite the open problem or point out the contradiction with the source
estimate.  Accordingly, this packet claims the identification, not novelty for
the monomial extremal itself.

\paragraph{Source question and estimate.}
On page 18, in Section~5 under ``Some open problems'', Anderson, Dritschel, and
Rovnyak ask whether there are Schur functions outside the Schur--Agler class on
the polydisk or unit ball which fail the estimates proved in Sections~3 and~4.
For the Euclidean ball $\B_d$, their Theorem~12, estimate~(42), says that every
ball Schur--Agler function $\varphi$ satisfies
\begin{equation}\label{ADR42}
 \left|\partial_1^{n_1}\cdots\partial_d^{n_d}\varphi(z)\right|
 \le d^{(n-1)/2}\,n_1!\cdots n_d!\,
 \frac{1-|\varphi(z)|^2}
 {(1-\|z\|_2^2)(1-\|z\|_2)^{n-1}},
 \qquad n=\sum_j n_j.
\end{equation}

\paragraph{Supporting literature.}
Dai--Chen--Pan, Theorem~2, proves that if
$f(z)=\sum_\alpha a_\alpha z^\alpha$ maps $\B_d$ to $\D$, then at the origin
\[
 |a_\alpha|\le
 \sqrt{\frac{|\alpha|^{|\alpha|}}{\alpha^\alpha}}
 (1-|a_0|^2),
\]
and describes the monomial extremals.  For the balanced multi-index
$\alpha=(1,\ldots,1)$, their sharp factor is $d^{d/2}$, while
\eqref{ADR42} permits only $d^{(d-1)/2}$.

\begin{theorem}
For every integer $d\ge2$, the polynomial
\[
 F_d(z)=d^{d/2}z_1z_2\cdots z_d
\]
belongs to the Schur class on $\B_d$, does not belong to the ball
Schur--Agler class, and violates estimate~\eqref{ADR42} at the origin by the
exact factor $\sqrt d$.  Hence Open Problem~2 has a full affirmative answer to
its existence question.
\end{theorem}

\begin{proof}
For $z\in\B_d$, arithmetic--geometric mean gives
\[
 |F_d(z)|
 =d^{d/2}\prod_{j=1}^d |z_j|
 \le d^{d/2}\left(\frac{\sum_j|z_j|^2}{d}\right)^{d/2}
 =\|z\|_2^d<1.
\]
Thus $F_d:\B_d\to\D$.  Its completely mixed derivative is constant:
\[
 \partial_1\partial_2\cdots\partial_dF_d(0)=d^{d/2}.
\]
In \eqref{ADR42}, take $n=d$, $n_1=\cdots=n_d=1$, and $z=0$.
Since $F_d(0)=0$, the proposed right side is exactly
$d^{(d-1)/2}$.  The quotient of the two sides is $\sqrt d>1$.
Therefore $F_d$ violates the estimate.  Theorem~12 proves the estimate for
every ball Schur--Agler function, so contraposition also shows that $F_d$ is
not Schur--Agler.
\end{proof}

\paragraph{Smallest instance.}
On $\B_2$, $F_2(z_1,z_2)=2z_1z_2$ is Schur because
$2|z_1z_2|\le |z_1|^2+|z_2|^2<1$.  At zero,
$|\partial_1\partial_2F_2|=2$, but (42) gives only $\sqrt2$.

\section*{The two examples singled out in the source}

The counterexample above settles existence, but it is different from the two
families the source explicitly asks about.  The following elementary companion
calculations settle both of those tests positively.

\begin{proposition}
Every homogeneous quadratic Schur polynomial on $\D^d$ satisfies all the
polydisk estimates of Anderson--Dritschel--Rovnyak.  In particular, the
Kaijser--Varopoulos polynomial satisfies them all.
\end{proposition}

\begin{proof}
Let $p:\D^d\to\D$ be homogeneous of degree two and put
$r=\|z\|_\infty$.  Homogeneity gives $|p(z)|\le r^2$.  The source's
first-derivative estimate is automatic from the one-variable Schwarz--Pick
lemma on the coordinate slice, because
\[
 |\partial_jp(z)|\le\frac{1-|p(z)|^2}{1-|z_j|^2}
 \le\frac{1-|p(z)|^2}
 {\sqrt{1-|z_j|^2}\sqrt{1-r^2}}.
\]
All derivatives of order at least three vanish.  If
$p=\sum_{|\alpha|=2}a_\alpha z^\alpha$, a torus Fourier integral gives
$|a_\alpha|\le1$.  Hence $|\partial_j^2p|\le2$, while
$|\partial_j\partial_kp|\le1$ for $j\ne k$.

For a second derivative, the right side of the source's permutation-sum
estimate is at least
\[
 2\frac{1-r^4}{1-r}\ge2.
\]
The right side of its multi-index estimate is at least
\[
 \alpha!\frac{1-r^4}{(1-r^2)(1-r)}
 =\alpha!\frac{1+r^2}{1-r}\ge\alpha!.
\]
These dominate the pure and mixed second derivatives, respectively.
\end{proof}

\begin{proposition}
For every Schur function $f:\B_d\to\D$, every pure-coordinate derivative
$\partial_j^nf$ satisfies both ball estimates in the source.  Consequently,
every Alpay--Kaptano\u{g}lu polynomial singled out in Open Problem~2 satisfies
all of those estimates.
\end{proposition}

\begin{proof}
Fix $z\in\B_d$ and a coordinate $j$.  Set
\[
 a=\|\widehat z_j\|_2,\qquad R=\sqrt{1-a^2},\qquad
 b=|z_j|,\qquad r=\|z\|_2.
\]
The slice $g(\lambda)=f(\widehat z_j+R\lambda e_j)$ maps $\D$ into
$\D$.  Applying the sharp one-variable higher Schwarz--Pick estimate at
$\lambda=z_j/R$ and rescaling gives
\begin{equation}\label{slice}
 |\partial_j^nf(z)|
 \le \frac{n!R(1-|f(z)|^2)}{(1-r^2)(R-b)^{n-1}}.
\end{equation}
But
\[
 R-b=\frac{1-r^2}{R+b}\ge\frac{1-r^2}{1+r}=1-r.
\]
Thus \eqref{slice} is at least as strong as the source's Theorem~10 bound for
a pure derivative, whose numerator contains $n!R$, and also as strong as its
Theorem~12 bound, whose numerator contains
$d^{(n-1)/2}n!\ge n!R$.  The first-order case is the standard infinitesimal
Schwarz--Pick inequality for the ball.

An Alpay--Kaptano\u{g}lu polynomial has the form
$z_1+\sum_{k=1}^m c_kz_2^{2k}$.  Every derivative of order at least two is
either zero or pure in $z_2$, while the first derivatives satisfy the universal
first-order estimate.  Hence all source estimates hold.
\end{proof}

\paragraph{Verification and provenance.}
The Schur bound and derivative comparison for $F_d$ are exact and use only
AM--GM.  The included checker confirms the exponent arithmetic for
$2\le d\le20$; it is not used as proof.  The run first searched its cheap
registry, solution, and attempt indexes for the exact arXiv id and the phrases
``Kaijser--Varopoulos'', ``Schur--Agler'', and ``Schwarz--Pick''.  Bounded web
and arXiv searches used the exact expressions ``$2z_1z_2$ Schur--Agler'',
``Kaijser--Varopoulos Schwarz--Pick derivative'', and the source title.  No
paper explicitly stating an answer to Open Problem~2 was found.  The closest
decisive source was arXiv:1109.2791, whose Theorem~2 supplies the general sharp
coefficient extremal but does not cite or identify the 2007 question.  This is
why the durable status is literature-implied rather than a new counterexample.

\paragraph{Scope.}
The existence question is completely answered on the unit ball in every
dimension $d\ge2$.  The packet does not produce a violating function on a
polydisk, settle sharpness of all source constants, or carry out the source's
Cartan-domain program.  The companion propositions settle only the two
explicit test families named in Open Problem~2.

\paragraph{Human-review recommendation.}
Verify the substitution $n=d$, $n_j=1$, $z=0$ in source estimate~(42), and
confirm the provenance classification against Dai--Chen--Pan Theorem~2.  The
mathematical counterexample itself is a one-line exact check.

\begin{thebibliography}{9}
\bibitem{ADR}
J. M. Anderson, M. A. Dritschel, and J. Rovnyak,
\emph{Schwarz--Pick inequalities for the Schur--Agler class on the polydisk
and unit ball}, Comput. Methods Funct. Theory 8 (2008), 339--361;
arXiv:math/0702269.

\bibitem{DCP}
S. Dai, H. Chen, and Y. Pan,
\emph{The Schwarz--Pick lemma of high order in several variables},
arXiv:1109.2791.

\bibitem{Rudin}
W. Rudin, \emph{Function Theory in the Unit Ball of $\mathbb C^n$},
Springer, 1980.
\end{thebibliography}

\end{document}

